\chapter{General Introduction}
\markboth{General Introduction}{General Introduction}

\section*{Introduction}

The rapid growth of cloud-native architectures over the past decade has profoundly
transformed how organisations deploy and manage their information systems. Modern
infrastructures, composed of dozens or even hundreds of servers, containers and
distributed services, generate operational data streams of unprecedented density.
In this context, system monitoring — once delegated to a network administrator
scrutinising a handful of graphs — has become a discipline in its own right, at
the heart of organisational resilience.

Yet current monitoring tools, however sophisticated, share a fundamental
limitation: they can \textit{detect} that a problem exists, but struggle to
explain \textit{why} it appeared. When faced with a ``CPU at 98\%'' alert, the
on-call engineer must still conduct a manual investigation lasting anywhere from
thirty minutes to several hours. This gap between \textit{detection} and
\textit{understanding} is precisely the problem this project sets out to solve.

\section{Context and Motivation}

Industry statistics confirm the scale of the challenge. According to PagerDuty's
\textit{State of Digital Operations} report (2023) \cite{pagerduty2023}, the
average cost of one hour of downtime exceeds \$300,000 for large enterprises,
while the mean time to resolution (\textsc{mttr}) exceeds 45 minutes in
organisations without mature \textsc{aiops} practices.

For small teams — a startup with two developers managing a \textsc{vps}, an
agency administering ten machines — the situation is even more critical. These
teams lack both the human resource of a dedicated DevOps engineer and the budget
for commercial solutions costing thousands of euros per month. They need a tool
that simply answers: ``\textit{Is my server healthy? If not, what should I do?}''

This is the context in which the \textbf{\caimp{}} project was born ---
\textit{Centralized AI Monitoring Platform} --- a self-hosted, open-source
platform that positions artificial intelligence not as a display gadget, but as
the core of the diagnostic process.

\section{Project Objectives}

\begin{enumerate}[label=\textbf{O\arabic*.}, leftmargin=2.2cm]
  \item \textbf{Answers-first dashboard}: design an interface whose home page
    immediately presents the global status, potential issues and recommended
    actions — requiring no monitoring expertise from the user.
  \item \textbf{Automatic change correlation}: develop an engine capable of
    linking each detected anomaly to recent environment modifications
    (deployments, restarts, package updates, SSH sessions) to automate root
    cause investigation.
  \item \textbf{Natural-language explanation}: integrate a local language model
    (without cloud API calls) to produce explanations understandable by a
    developer with no DevOps training.
  \item \textbf{Resource saturation forecasting}: implement a Holt-Winters
    forecasting engine to anticipate threshold breaches over the next 24 hours.
  \item \textbf{Robust microservices architecture}: design a decoupled architecture
    based on an event bus (NATS JetStream) supporting horizontal scalability and
    multi-tenant data isolation.
  \item \textbf{End-to-end security}: implement mTLS for agents, JWT management,
    and RLS-based data isolation per organisation.
  \item \textbf{Accessible deployment}: make full installation achievable in
    under ten minutes via Docker Compose, on any machine with 6\,GB of RAM.
\end{enumerate}

\section{Working Methodology}

The development of \caimp{} followed an agile approach adapted to the academic
context. The three-person team held daily synchronisation sessions
(\textit{stand-ups}) and organised its work in weekly sprints of five days.
The code repository was hosted on GitHub, with a protected main branch and a
development branch from which features were developed and tested before
integration. Commit messages followed the \textit{Conventional Commits}
convention.

\section{Report Structure}

\begin{itemize}
  \item \textbf{Chapter 1}: General context, modern observability, and the AIOps domain.
  \item \textbf{Chapter 2}: State of the art of existing monitoring solutions.
  \item \textbf{Chapter 3}: Functional and non-functional requirements, UML diagrams.
  \item \textbf{Chapter 4}: Platform architecture, data model, and processing flows.
  \item \textbf{Chapter 5}: Technology choices and justifications.
  \item \textbf{Chapter 6}: Implementation details with annotated code extracts.
  \item \textbf{Chapter 7}: Testing strategy and validation results.
  \item \textbf{Chapter 8}: Deployment and production procedures.
  \item \textbf{Chapter 9}: Results, discussion, and perspectives.
  \item \textbf{Chapter 10}: General conclusion.
\end{itemize}

\section*{Conclusion}

This project represents the culmination of one month of intensive work, during
which a complete infrastructure monitoring platform was designed, developed,
tested and documented. The following chapters demonstrate the scientific and
technical rigour we sought to maintain throughout this process. The next
chapter establishes the contextual foundations essential to understanding the
work as a whole.
